% =====================================================================
\section{Thảo luận}
\label{sec:discussion}
% =====================================================================

\subsection{Trả lời ba câu hỏi đánh giá}
\label{sec:disc-answers}

\paragraph{C1 --- Hiệu quả phân loại.}
\RFindingCone

\paragraph{C2 --- Vai trò của dữ liệu cho nhiệm vụ phụ.}
\RFindingCtwo

\paragraph{C3 --- Chi phí thực hiện.}
\RFindingCthree

\subsection{Đọc kết quả theo hướng xu hướng ngân sách nhãn}
\label{sec:disc-trend}

Giả thuyết ban đầu cho rằng nhiệm vụ phụ có lợi nhất khi ít nhãn, vì lúc đó tín hiệu tự giám
sát bổ sung được thông tin mà nhãn có giám sát không cung cấp đủ. Hình~\ref{fig:macro-f1} cho
phép kiểm tra giả thuyết này bằng cách so sánh độ lớn của $\Delta$ ở ba mức $10\,\%$, $20\,\%$
và $50\,\%$. Điều cần chú ý là \emph{độ lớn của $\Delta$}, không chỉ dấu của nó: một cải thiện
$+0{,}3$ điểm phần trăm nằm trong cùng bậc độ lớn với độ lệch chuẩn quan sát được và do đó
không thể phân biệt với nhiễu.

\subsection{Các yếu tố gây nhiễu và mối đe dọa đến tính hợp lệ}
\label{sec:disc-threats}

\paragraph{Nhiễu do thống kê BatchNorm.}
\label{sec:exempt-bn}
Như đã nêu ở \secref{sec:arch}, encoder dùng chung các lớp BatchNorm cho cả hai nhánh. Khi so
sánh \cfg{E2-L} với \cfg{E1}, hai cấu hình khác nhau ở cả hàm mục tiêu \emph{lẫn} thành phần
của batch đi qua BatchNorm. Vì vậy $\Delta_s$ của phép so sánh này đo \emph{tổng hợp} tác dụng
của nhiệm vụ phụ và thay đổi chuẩn hóa, không tách riêng được hai thành phần. Đây là hệ quả của
một quyết định thiết kế kiến trúc (dùng chung BatchNorm là cần thiết để giảm số tham số và giữ
so sánh tham số ở mức ngang bằng), không phải lỗi cài đặt.

\paragraph{Nhiễu do chọn checkpoint bằng validation.}
Mọi cấu hình đều được chọn checkpoint bằng Macro-F1 validation cao nhất. Vì chỉ có 5.000 ảnh
validation, tiêu chí này có phương sai riêng: chênh lệch nhỏ về Macro-F1 test có thể chỉ phản
ánh việc \enquote{may mắn} khi chọn epoch. Việc dùng cùng tiêu chí cho mọi cấu hình làm giảm
nhưng không loại bỏ được vấn đề này.

\paragraph{Chỉ ba lần lặp.}
Ba lần lặp cho phép ước lượng một độ lệch chuẩn mẫu nhưng khoảng tin cậy vẫn rất rộng. Với
$n=3$, một chênh lệch trung bình bằng $0$ hoàn toàn tương thích với các lần lặp có $\Delta_s$
trái dấu. Do đó báo cáo này chỉ trình bày \emph{hướng} của hiệu ứng.

\paragraph{Phép chia dữ liệu cố định.}
Cả ba lần lặp dùng chung một tập validation và một tập test. Phương sai do chính phép chia gây
ra \textbf{không} được ước lượng. Muốn ước lượng thành phần này cần lặp lại cả phép chia, tức
tăng chi phí lên đáng kể.

\paragraph{Không dùng Early Stopping.}
Quyết định bỏ Early Stopping để bảo đảm số bước cập nhật có giám sát giống nhau. Đánh đổi là
một số lượt chạy có thể đã quá khớp (xem loss huấn luyện rất nhỏ ở cuối giai đoạn, ví dụ
$L_{\mathrm{sup}} < 0{,}01$ ở epoch 100 tại mức $10\,\%$ nhãn), và checkpoint được chọn gần
cuối quá trình. Đây là hành vi nhất quán giữa các cấu hình, nhưng số liệu cuối cùng không phản
ánh chất lượng tốt nhất mà cấu hình đó có thể đạt được dưới một chiến lược dừng khác.

\paragraph{Kích thước batch và chất lượng BatchNorm.}
Với $|\DL| = 4.500$ ở mức $10\,\%$, một epoch chỉ có 36 bước cập nhật. Số bước nhỏ làm cho
thống kê BatchNorm kém ổn định hơn so với trường hợp nhiều dữ liệu. Đây là hiệu ứng có thật của
thiết lập ít nhãn và không nên che giấu.

\subsection{Những gì đã được cưỡng chế bằng mã nguồn, không chỉ bằng văn bản}
\label{sec:disc-enforced}

Một phần giá trị của báo cáo này nằm ở chỗ các ràng buộc về phương pháp được \emph{cưỡng chế}
bằng mã chứ không chỉ được ghi trong văn bản:

\begin{itemize}[leftmargin=1.2cm,itemsep=3pt]
  \item Hàm \code{make\_run\_loaders} nhận tham số \code{rotation\_source}; với \cfg{E1} giá
        trị là \code{"none"} và \textbf{không có} DataLoader nào cho nhánh xoay được tạo ra.
        Do đó không thể vô tình đưa ảnh xoay qua BatchNorm trong cấu hình baseline.
  \item Dataset \code{RotationViewDataset} là cấu trúc duy nhất biết $\DU$; nó chứa chỉ mục
        nhưng \textbf{không} chứa nhãn siêu lớp, nên một lỗi lập trình ở nhánh phụ không thể
        vô tình đọc nhãn thật.
  \item Bộ sinh số ngẫu nhiên riêng cho nhánh phụ được khởi tạo tách khỏi RNG chính, nên thêm
        nhánh phụ không làm thay đổi augmentation của nhánh phân loại.
  \item Tập test chỉ được truy cập ở hàm \code{evaluate} cuối cùng sau khi checkpoint đã được
        cố định bằng validation.
\end{itemize}

\subsection{Vì sao contrastive learning không tự động giúp phân loại}
\label{sec:disc-contrastive}

Khối mở rộng \cfg{E3}--\cfg{E6} cung cấp một đối chứng quan trọng cho trực giác thường gặp rằng
\enquote{thêm học tự giám sát thì tốt hơn}. Có ít nhất bốn lý do khiến nhiệm vụ phụ contrastive
không đương nhiên cải thiện nhiệm vụ phân loại siêu lớp:

\begin{enumerate}[leftmargin=1.2cm,itemsep=3pt]
  \item \textbf{Mục tiêu khác nhau.} NT-Xent học bất biến trước \emph{tăng cường ảnh} (cắt, đổi
        màu, chuyển xám). Animal và Vehicle phân biệt nhau bằng \emph{hình dạng ngữ nghĩa}, mà
        tăng cường mạnh lại chính là thứ \emph{phá} thông tin cấp thấp hữu ích. Hai mục tiêu
        không mâu thuẫn về nguyên tắc nhưng không trùng nhau.
  \item \textbf{Encoder bị chia sẻ với hai mục tiêu cạnh tranh.} Vì $\enc$ là duy nhất cho mọi
        nhánh, gradient của loss contrastive có thể kéo encoder về vùng làm giảm chất lượng đặc
        trưng cho phân loại ngay ở giai đoạn đầu — khi đó gradient phân loại còn yếu vì chỉ có
        $|\DL|$ nhỏ.
  \item \textbf{Chi phí phải trả bằng chính ngân sách cập nhật.} Mỗi bước vẫn chỉ thực hiện một
        bước optimizer. Việc thêm hai view contrastive làm thay đổi phân bố gradient của bước
        đó, chứ không \enquote{thêm} bước học. Đây chính là lý do vì sao so sánh phải \emph{ghép
        cặp theo cùng số bước cập nhật}, và vì sao không thể diễn giải kết quả mở rộng như
        \enquote{cùng ngân sách tính toán}.
  \item \textbf{Nguồn ảnh không nhãn không tự động chuyển thành tín hiệu hữu ích.} E3/E5 dùng
        cả $\Dtrain$ (gấp 10 lần $\DL$ ở mức $10\,\%$ nhãn), nhưng ảnh không nhãn chỉ cung cấp
        tín hiệu nếu \emph{bất biến mà chúng dạy được} trùng với bất biến mà nhiệm vụ chính cần.
\end{enumerate}

\RExtensionAnalysis

Một điểm phương pháp luận đáng nhấn mạnh: $\lambda_c$ của báo cáo này được chọn bằng
Macro-F1 \emph{validation} với một script không đánh giá test
(\secref{sec:results-ext-lambda}). Nếu cho phép chọn $\lambda_c$ bằng test, kết quả của khối mở
rộng sẽ trông tốt hơn một cách giả tạo. Đây là loại sai lệch mà đề cương yêu cầu phải tránh.

\subsection{Hướng cải thiện}
\label{sec:disc-next}

Theo thứ tự ưu tiên dựa trên kết quả thu được:

\begin{enumerate}[leftmargin=1.2cm,itemsep=3pt]
  \item \textbf{Tăng số lần lặp} lên tối thiểu $5$--$10$ ở mức $10\,\%$ nhãn, mức được dự đoán
        nhạy nhất với nhiệm vụ phụ. Đây là thay đổi cho tỷ lệ tín hiệu/nhiễu tốt nhất trên mỗi
        đơn vị chi phí.
  \item \textbf{Tách BatchNorm theo nhánh} (ví dụ dùng \code{Affine} riêng cho nhánh phụ thay vì
        dùng chung thống kê chạy) để có thể quy kết hiệu ứng cho nhiệm vụ phụ thay vì cho thay
        đổi chuẩn hóa.
  \item \textbf{Khảo sát $\lamrot$} trên lưới đã khai báo $\{0{,}1;\ 0{,}25;\ 0{,}5;\ 1{,}0\}$
        ở mức $10\,\%$ nhãn, kèm kiểm tra Rotation Accuracy để biết nhánh phụ có thực sự được
        học hay không.
  \item \textbf{Thay đổi tỷ lệ đóng góp} giữa hai nhánh bằng cách chỉnh kích thước $B_R$ thay
        vì chỉ chỉnh $\lamrot$.
  \item \textbf{Chỉ sau khi phần chính đã đủ} mới triển khai \emph{một} nhánh mở rộng contrastive
        (\tabref{tab:extensions}), ưu tiên SupCon vì nó chỉ dùng $\DL$ và do đó không làm phức
        tạp thêm câu hỏi về việc che nhãn.
\end{enumerate}
